% Seccion: Energía solar térmica
\section{Energía solar térmica}
\label{sec:cap02_energia_solar}

La energía solar térmica aprovecha la radiación electromagnética proveniente del sol para calentar fluidos de trabajo.
La constante solar en el límite exterior de la atmósfera terrestre posee un valor promedio de mil trescientos sesenta y siete vatios por metro cuadrado.
Al atravesar la atmósfera, la radiación experimenta atenuaciones por absorción y dispersión de gases y aerosoles.
La radiación global sobre una superficie inclinada se compone de la fracción directa, difusa y la reflejada por el suelo o albedo.
Para maximizar la captación en latitudes medias, es fundamental inclinar el plano del colector de acuerdo con la latitud local y la estación del año.
